\chapter{{\fontsize{14pt}{21pt}\selectfont\MakeUppercase{Техническая реализация системы}}}
\label{cha:design}

\section{Разработка пользовательского интерфейса приложения}

Для реализации пользовательского интерфейса веб-приложения FoodPlanner использовались React и библиотека компонентов Mantine. Выбранный технологический стек позволил сократить время разработки, обеспечить единый визуальный стиль экранов и повысить предсказуемость поведения интерфейса.

Главная страница выполняет функцию стартового навигационного узла и предоставляет пользователю структурированное представление о возможностях платформы. На странице размещены ключевые переходы к тематическим разделам системы: рецептам, планам питания, конструктору, отзывам и личному кабинету. Такое решение обеспечивает быстрое вхождение в систему без необходимости предварительного изучения сложной структуры меню.

Композиционно страница организована как набор контентных секций и карточек предпросмотра. Карточный формат визуально разделяет сущности и помогает пользователю быстро определить целевой сценарий. С точки зрения UX это снижает время поиска нужного действия и повышает вероятность перехода в функциональные разделы. Внешний вид главной страницы приложения представлен на рисунке~\ref{fig:image_7}.

\begin{figure}[H]
  \centering
  \IfFileExists{images/image_7.png}{
    \includegraphics[width=0.95\textwidth]{image_7.png}
  }{
    \fbox{\parbox{0.9\textwidth}{\centering Место для изображения главной страницы\\
    Файл: \texttt{Report/images/image_7.png}}}
  }
  \caption{Главная страница приложения}
  \label{fig:image_7}
\end{figure}

Раздел рецептов реализует каталог с возможностью фильтрации по ключевым параметрам питания и кулинарного процесса. Пользователь может формировать релевантную выборку рецептов по тегам, калорийности, времени приготовления и другим критериям. Использование серверной фильтрации и постраничной загрузки поддерживает высокую отзывчивость интерфейса при росте объема данных.

Дополнительно реализована логика корректного отображения элементов пагинации: кнопка <<Показать еще>> отображается только при наличии данных для следующей подгрузки. Карточки рецептов содержат ключевые атрибуты для первичного выбора: название, краткое описание, теги, время приготовления и базовые показатели. Интерфейс раздела рецептов представлен на рисунке~\ref{fig:image_8}.

\begin{figure}[H]
  \centering
  \IfFileExists{images/image_8.png}{
    \includegraphics[width=0.95\textwidth]{image_8.png}
  }{
    \fbox{\parbox{0.9\textwidth}{\centering Место для изображения страницы рецептов\\
    Файл: \texttt{Report/images/image_8.png}}}
  }
  \caption{Страница рецептов}
  \label{fig:image_8}
\end{figure}

Детальная страница рецепта объединяет информационный и прикладной контуры. Информационный блок включает название, описание, пищевую ценность (калории, белки, жиры, углеводы), теги и пошаговую инструкцию. Прикладной блок реализует интерактивные действия: управление ингредиентами, формирование списка покупок, сохранение в профиль и экспорт.

Особое значение имеет блок ингредиентов с чекбоксами: пользователь отмечает наличие позиций, после чего система формирует список недостающих покупок. При повторном посещении страницы состояние чекбоксов восстанавливается, что поддерживает непрерывность пользовательского сценария. Детальная страница рецепта представлена на рисунке~\ref{fig:image_9}.

\begin{figure}[H]
  \centering
  \IfFileExists{images/image_9.png}{
    \includegraphics[width=0.95\textwidth]{image_9.png}
  }{
    \fbox{\parbox{0.9\textwidth}{\centering Место для изображения детальной страницы рецепта\\
    Файл: \texttt{Report/images/image_9.png}}}
  }
  \caption{Страница конкретного рецепта}
  \label{fig:image_9}
\end{figure}

Раздел планов питания является витриной готовых рационов и ориентирован на задачу быстрого выбора стратегии питания. Пользователь получает список планов с фильтрацией по типу, калорийности, рейтингу и дополнительным параметрам, включая персонализированные предпочтения. Серверная пагинация позволяет обрабатывать большие массивы данных без перегрузки интерфейса.

Карточка плана содержит ключевые характеристики для принятия решения: автор, тип плана, длительность, агрегированные показатели и оценка. Унификация карточек между разделами сокращает порог обучения и обеспечивает консистентный пользовательский опыт. Интерфейс страницы планов питания представлен на рисунке~\ref{fig:image_10}.

\begin{figure}[H]
  \centering
  \IfFileExists{images/image_10.png}{
    \includegraphics[width=0.95\textwidth]{image_10.png}
  }{
    \fbox{\parbox{0.9\textwidth}{\centering Место для изображения страницы планов питания\\
    Файл: \texttt{Report/images/image_10.png}}}
  }
  \caption{Страница планов питания}
  \label{fig:image_10}
\end{figure}

Детальная страница плана питания раскрывает структуру рациона по дням и приемам пищи (завтрак, обед, ужин, перекусы) с привязкой к рецептам и пищевым параметрам. Такой формат позволяет оценить рацион не только по агрегированным показателям, но и на уровне практического применения по каждому дню.

Интерфейс поддерживает раскрывающиеся секции и группировку данных, что упрощает навигацию внутри длинных планов и снижает визуальный шум. Как и на странице рецепта, доступны действия прикладного уровня: работа с ингредиентами, формирование списка покупок и сохранение результатов. Детальная страница плана питания представлена на рисунке~\ref{fig:image_11}.

\begin{figure}[H]
  \centering
  \IfFileExists{images/image_11.png}{
    \includegraphics[width=0.95\textwidth]{image_11.png}
  }{
    \fbox{\parbox{0.9\textwidth}{\centering Место для изображения детальной страницы плана\\
    Файл: \texttt{Report/images/image_11.png}}}
  }
  \caption{Страница конкретного плана питания}
  \label{fig:image_11}
\end{figure}

Конструктор является центральным функциональным модулем приложения и предназначен для создания персонализированных планов питания. Интерфейс позволяет задавать параметры плана, выбирать рецепты из каталога и распределять их по календарной структуре. Реализована визуальная сегментация дней и поддержка вариативного количества перекусов.

Пользователь одновременно видит доступный контент и текущую структуру плана, что упрощает составление и корректировку рациона. По завершении формирования план отправляется на модерацию, после чего рабочая область конструктора подготавливается к следующему сценарию. Интерфейс конструктора представлен на рисунке~\ref{fig:image_12}.

\begin{figure}[H]
  \centering
  \IfFileExists{images/image_12.png}{
    \includegraphics[width=0.95\textwidth]{image_12.png}
  }{
    \fbox{\parbox{0.9\textwidth}{\centering Место для изображения конструктора плана\\
    Файл: \texttt{Report/images/image_12.png}}}
  }
  \caption{Страница конструктора плана питания}
  \label{fig:image_12}
\end{figure}

Кабинет модератора реализован как специализированный интерфейс принятия решений. Для модератора скрыты нерелевантные персональные блоки, а в фокус вынесены заявки на проверку. Это снижает информационный шум и поддерживает операционный характер работы.

В блоке модерации рецептов отображаются карточки материалов в статусе проверки с ключевыми метаданными: тип сущности, автор, дата подачи и название. Каждая карточка предоставляет действия одобрения/отклонения и переход к детальному просмотру. Интерфейс кабинета модератора представлен на рисунке~\ref{fig:image_13}.

\begin{figure}[H]
  \centering
  \IfFileExists{images/image_13.png}{
    \includegraphics[width=0.95\textwidth]{image_13.png}
  }{
    \fbox{\parbox{0.9\textwidth}{\centering Место для изображения кабинета модератора\\
    Файл: \texttt{Report/images/image_13.png}}}
  }
  \caption{Личный кабинет модератора с модерацией}
  \label{fig:image_13}
\end{figure}

Личный кабинет пользователя предназначен для управления персональными данными и параметрами питания, влияющими на персонализацию контента. Пользователь может обновлять профиль, указывать цели, предпочтительную диету, особенности питания и любимые теги.

С точки зрения UX профиль построен как набор самостоятельных функциональных блоков: данные аккаунта, персональные предпочтения, избранное, собственные материалы, статусы модерации и связанные сущности (включая списки покупок). Поддержка пустых состояний и уведомлений о результате операций обеспечивает понятную обратную связь пользователю. Внешний вид личного кабинета пользователя представлен на рисунках~\ref{fig:image_14} и~\ref{fig:image_15}.

\begin{figure}[H]
  \centering
  \IfFileExists{images/image_14.png}{
    \includegraphics[width=0.95\textwidth]{image_14.png}
  }{
    \fbox{\parbox{0.9\textwidth}{\centering Место для изображения личного кабинета (вариант 1)\\
    Файл: \texttt{Report/images/image_14.png}}}
  }
  \caption{Личный кабинет пользователя}
  \label{fig:image_14}
\end{figure}

\begin{figure}[H]
  \centering
  \IfFileExists{images/image_15.png}{
    \includegraphics[width=0.95\textwidth]{image_15.png}
  }{
    \fbox{\parbox{0.9\textwidth}{\centering Место для изображения личного кабинета (вариант 2)\\
    Файл: \texttt{Report/images/image_15.png}}}
  }
  \caption{Личный кабинет пользователя}
  \label{fig:image_15}
\end{figure}

ИИ-ассистент интегрирован в приложение как доступный с любого экрана диалоговый интерфейс. Пользователь может открыть чат через плавающую кнопку и получить консультацию по рациону, подбору блюд и общим вопросам питания. Реализация в формате выдвижной панели обеспечивает быстрый доступ к функции без прерывания основного пользовательского сценария.

Интерфейс чата включает историю сообщений, поле ввода с повышенной читаемостью и явные состояния отправки/ожидания ответа. Дополнительно реализовано форматирование длинных ответов и обработка системных ошибок внешнего API через уведомления. Открытое диалоговое окно с ИИ-ассистентом представлено на рисунке~\ref{fig:image_16}.

\begin{figure}[H]
  \centering
  \IfFileExists{images/image_16.png}{
    \includegraphics[width=0.95\textwidth]{image_16.png}
  }{
    \fbox{\parbox{0.9\textwidth}{\centering Место для изображения окна ИИ-ассистента\\
    Файл: \texttt{Report/images/image_16.png}}}
  }
  \caption{Открытое диалоговое окно с ИИ-ассистентом}
  \label{fig:image_16}
\end{figure}

Раздел оценок и отзывов обеспечивает прозрачный механизм социальной валидации контента. Пользователь может анализировать оценки и текстовые комментарии других участников платформы, что повышает доверие к опубликованным материалам и помогает принимать обоснованные решения при выборе рецептов и планов питания.

Функционально раздел агрегирует пользовательский опыт и одновременно является каналом обратной связи для авторов контента и модераторов. Интерфейс экрана реализован в формате карточек и списков с акцентом на читаемость комментариев, рейтинг и связь с объектом оценки. Страница оценок и отзывов представлена на рисунке~\ref{fig:image_17}.

\begin{figure}[H]
  \centering
  \IfFileExists{images/image_17.png}{
    \includegraphics[width=0.95\textwidth]{image_17.png}
  }{
    \fbox{\parbox{0.9\textwidth}{\centering Место для изображения страницы оценок и отзывов\\
    Файл: \texttt{Report/images/image_17.png}}}
  }
  \caption{Страница оценок и отзывов}
  \label{fig:image_17}
\end{figure}

\section{Техническая реализация серверной части}

Информация о выборе архитектурного подхода и программных средств приведена в разделе~1.7. В данном разделе рассматриваются практические аспекты непосредственной серверной разработки: структура проекта, организация слоев ответственности, работа с базой данных, реализация API и поддержка ключевых бизнес-процессов.

Серверная часть FoodPlanner реализована на языке Python с использованием Django и Django REST Framework. Такой выбор позволил реализовать прикладную логику в сжатые сроки и при этом сохранить высокий уровень структурированности кода. В качестве основной СУБД используется PostgreSQL; для локальной отладки и быстрых проверок может использоваться SQLite.

Структура директории проекта серверной части. Серверный код размещен в директории \texttt{Backend} и организован по доменным приложениям. Каждое приложение изолирует отдельную часть предметной области и имеет собственные модели, сериализаторы, представления и маршруты. Ниже приведено назначение ключевых директорий:
\begin{compactlist}
  \item \texttt{/Backend} --- корневая директория серверной части, содержащая конфигурацию проекта, доменные модули, миграции и инфраструктурные файлы;
  \item \texttt{/Backend/config} --- модуль глобальной конфигурации: параметры Django, middleware, CORS/CSRF, подключение приложений, верхнеуровневые URL и служебные endpoint'ы;
  \item \texttt{/Backend/users} --- модуль пользователей и ролей: регистрация, вход/выход, профиль, ролевой доступ;
  \item \texttt{/Backend/recipes} --- модуль рецептов и справочников: категории, теги, ингредиенты, рецепты, фильтрация и модерационные действия;
  \item \texttt{/Backend/meal\_plans} --- модуль планов питания и shopping list: создание, редактирование, отправка на модерацию, детализация по дням и приемам пищи;
  \item \texttt{/Backend/interactions} --- модуль пользовательских взаимодействий: отзывы, оценки, жалобы, избранное;
  \item \texttt{/Backend/notifications} --- модуль уведомлений по событиям платформы;
  \item \texttt{/Backend/moderation} --- модуль журналирования модерационных действий;
  \item \texttt{/Backend/assistant} --- модуль ИИ-ассистента: серверный прокси-запрос к внешнему LLM-провайдеру с обработкой ошибок и контекста;
  \item \texttt{/Backend/*/migrations} --- миграции схемы БД для каждого доменного приложения;
  \item \texttt{/Backend/users/management/commands} --- management-команды (в том числе \texttt{seed\_demo\_data}) для ускорения тестирования сценариев.
\end{compactlist}

Организация серверной части по указанным модулям обеспечивает разделение ответственности, упрощает тестирование и снижает стоимость последующих изменений. Структура серверной части представлена на рисунке~\ref{fig:image_18}.

\begin{figure}[H]
  \centering
  \IfFileExists{images/image_18.png}{
    \includegraphics[width=0.95\textwidth]{image_18.png}
  }{
    \fbox{\parbox{0.9\textwidth}{\centering Место для изображения структуры серверной части\\
    Файл: \texttt{Report/images/image_18.png}}}
  }
  \caption{Структура серверной части проекта}
  \label{fig:image_18}
\end{figure}

Слойная организация серверной логики. Серверная часть построена по принципу разделения ответственности между тремя уровнями:
\begin{compactlist}
  \item уровень API (views/viewsets + serializers) --- прием входящих запросов, валидация формата данных, маршрутизация, формирование ответов;
  \item уровень бизнес-логики (доменные правила в приложениях, permissions, actions) --- реализация прикладных сценариев, проверка прав доступа, обработка статусов модерации и переходов жизненного цикла сущностей;
  \item уровень данных (models + ORM + migrations) --- описание структуры БД, связей, ограничений целостности, индексов и миграций.
\end{compactlist}

Такая модель не допускает концентрации всей логики в одном слое: API-уровень отвечает за протокол взаимодействия, доменный уровень --- за поведение системы, а модельный уровень --- за целостность хранения данных.

Реализация API и ключевые серверные сценарии. Сервер предоставляет REST API для всех основных ролей и пользовательских сценариев платформы. Ключевые контуры API включают:
\begin{compactlist}
  \item пользователей и профиль (регистрация, авторизация, выход, чтение и обновление профиля);
  \item рецепты (каталог, фильтрация, поиск, детальная карточка, модерация);
  \item планы питания (создание, редактирование, отправка на модерацию, решение модератора);
  \item отзывы и оценки (создание, изменение, модерация);
  \item жалобы и избранное (подача и обработка жалоб, пользовательские коллекции);
  \item списки покупок (формирование из рецепта/плана и сохранение в профиль);
  \item уведомления по событиям модерации и пользовательским действиям;
  \item ИИ-ассистент через безопасный backend-proxy endpoint.
\end{compactlist}

Работа с базой данных и моделью данных. В рассматриваемом проекте взаимодействие с БД реализовано преимущественно через Django ORM. В отличие от подхода с полностью ручным формированием SQL-запросов, ORM-подход был выбран осознанно и обусловлен следующими факторами:
\begin{compactlist}
  \item повышение скорости и надежности разработки за счет типовых операций чтения/записи без дублирования SQL-кода;
  \item единообразие и предсказуемость доступа к данным во всех доменных модулях проекта;
  \item снижение вероятности ошибок в CRUD-операциях и упрощение сопровождения при изменении схемы.
\end{compactlist}

Использование ORM в проекте не исключает контроля производительности. Для ресурсоемких сценариев применяются индексы, ограничения целостности, оптимизация выборок через механизмы ORM и явное проектирование связей между сущностями. Такая комбинация позволяет сохранить баланс между удобством разработки и эффективностью выполнения запросов.

Модель данных включает сущности пользователей, ролей, рецептов, ингредиентов, планов питания, элементов плана, отзывов, жалоб, избранного, уведомлений, списков покупок и журналов модерации. Между сущностями используются связи <<один-ко-многим>> и <<многие-ко-многим>>, что позволяет корректно покрывать прикладные сценарии.

Для критичных направлений (поиск, модерация, привязка shopping list к источнику, контроль уникальности взаимодействий) введены ограничения целостности и индексы. Миграции фиксируют эволюцию схемы БД и обеспечивают воспроизводимое развертывание на разных окружениях.

Хранение прикладных данных и медиа. Основные транзакционные данные (профили, рецепты, планы, отзывы, жалобы, уведомления) сохраняются в PostgreSQL. Файловые ресурсы в текущей версии могут размещаться в файловой системе сервера, что достаточно для учебного и пилотного контура эксплуатации. При переходе к промышленной эксплуатации архитектура допускает перенос медиа-ресурсов в объектное S3-совместимое хранилище без изменения прикладной бизнес-логики, за счет абстракции слоя хранения.

Пример получения пользователя при авторизации через сериализатор DRF приведен в листинге~\ref{listing:user-auth-lookup}. В отличие от ручного сравнения паролей в SQL-запросе, в проекте используется встроенный механизм \texttt{authenticate}, который учитывает хеширование паролей и backend-аутентификации Django.

\begin{listing}[H]
\caption{Проверка учетных данных пользователя при входе}
\label{listing:user-auth-lookup}
\begin{mdframed}
\begin{minted}[breaklines=true]{python}
class LoginSerializer(serializers.Serializer):
    email = serializers.EmailField()
    password = serializers.CharField(write_only=True)

    def validate(self, attrs):
        user = authenticate(
            request=self.context.get("request"),
            username=attrs["email"],
            password=attrs["password"],
        )
        if not user:
            raise serializers.ValidationError("Invalid credentials.")
        attrs["user"] = user
        return attrs
\end{minted}
\end{mdframed}
\end{listing}

Конфигурация подключения к базе данных определяется централизованно через настройки приложения и переменные окружения (листинг~\ref{listing:db-settings-singleton-analog}). Такой подход выполняет ту же функцию, что и ручная singleton-инициализация подключения в других стеках, но в терминах Django и его жизненного цикла.

\begin{listing}[H]
\caption{Конфигурация подключения к PostgreSQL в Django}
\label{listing:db-settings-singleton-analog}
\begin{mdframed}
\begin{minted}[breaklines=true]{python}
if env("POSTGRES_DB", default=""):
    DATABASES = {
        "default": {
            "ENGINE": "django.db.backends.postgresql",
            "NAME": env("POSTGRES_DB"),
            "USER": env("POSTGRES_USER", default="postgres"),
            "PASSWORD": env("POSTGRES_PASSWORD", default="postgres"),
            "HOST": env("POSTGRES_HOST", default="localhost"),
            "PORT": env("POSTGRES_PORT", default="5432"),
        }
    }
else:
    DATABASES = {
        "default": {
            "ENGINE": "django.db.backends.sqlite3",
            "NAME": BASE_DIR / "db.sqlite3",
        }
    }
\end{minted}
\end{mdframed}
\end{listing}

Использование переменных окружения является обязательной частью конфигурации backend. Переменные хранят параметры подключения, секретные ключи и настройки окружения без их жесткого встраивания в исходный код. Это повышает безопасность и позволяет использовать различные значения для разработки и production-среды (листинг~\ref{listing:env-loading}).

\begin{listing}[H]
\caption{Загрузка переменных окружения в настройках backend}
\label{listing:env-loading}
\begin{mdframed}
\begin{minted}[breaklines=true]{python}
import environ

BASE_DIR = Path(__file__).resolve().parent.parent
env = environ.Env(
    DEBUG=(bool, True),
    SECRET_KEY=(str, "django-insecure-change-me"),
)
environ.Env.read_env(BASE_DIR / ".env")

DEBUG = env("DEBUG")
SECRET_KEY = env("SECRET_KEY")
\end{minted}
\end{mdframed}
\end{listing}

Персонализация и серверная фильтрация. Одной из ключевых задач backend является выдача релевантных рецептов и планов питания с учетом параметров профиля (цели, теги, особенности питания). Фильтрация и пагинация реализованы на сервере, что обеспечивает:
\begin{compactlist}
  \item корректную обработку больших выборок;
  \item снижение нагрузки на клиент;
  \item единый источник правды по составу результатов;
  \item стабильное поведение интерфейса при комбинировании фильтров.
\end{compactlist}

Модерация как самостоятельный процесс. Модерационный контур реализован отдельно от пользовательского потока: пользователь создает контент и отправляет его на проверку, модератор или администратор принимает решение, после чего статус сущности обновляется и пользователю направляется уведомление. Такой подход обеспечивает управляемый жизненный цикл контента, аудит действий и масштабируемость процесса при росте нагрузки.

Используемые внешние библиотеки. Практическая реализация серверной части опирается на следующие группы библиотек:
\begin{compactlist}
  \item Django/DRF --- веб-фреймворк и REST-интерфейс;
  \item psycopg --- драйвер подключения к PostgreSQL;
  \item drf-spectacular --- генерация OpenAPI/Swagger-схемы;
  \item django-cors-headers --- управление CORS-политиками;
  \item django-environ --- загрузка конфигурации из переменных окружения;
  \item gunicorn --- production WSGI-сервер;
  \item стандартные модули Python для сетевого взаимодействия с внешним API ИИ-ассистента.
\end{compactlist}

Описание конечных точек API. Взаимодействие между клиентской и серверной частями реализовано в стиле REST с использованием протокола HTTP. Обмен данными строится на стандартных методах:
\begin{compactlist}
  \item \texttt{GET} --- получение ресурса или коллекции ресурсов;
  \item \texttt{POST} --- создание ресурса или запуск прикладного действия;
  \item \texttt{PUT/PATCH} --- полное или частичное обновление существующего ресурса;
  \item \texttt{DELETE} --- удаление ресурса (или логическое удаление в рамках бизнес-логики).
\end{compactlist}

Каждой группе ресурсов соответствует набор endpoint'ов, объединенных общим префиксом \texttt{/api}. Для отладки и тестирования API на этапе разработки использовались Swagger-интерфейс (\texttt{/api/docs/}) и ручная проверка запросов в Insomnia/Postman.

Список ключевых endpoint'ов серверной части:
\begin{compactlist}
  \item аутентификация и профиль пользователя:
  \texttt{POST /api/users/register/},
  \texttt{POST /api/users/login/},
  \texttt{POST /api/users/logout/},
  \texttt{GET/PATCH /api/users/me/};
  \item администрирование пользователей:
  \texttt{GET/POST /api/users/manage/},
  \texttt{GET/PATCH/DELETE /api/users/manage/\{id\}/};
  \item рецепты и справочники:
  \texttt{GET/POST /api/recipes/},
  \texttt{GET/PATCH/DELETE /api/recipes/\{id\}/},
  \texttt{POST /api/recipes/\{id\}/moderate/},
  \texttt{GET /api/recipes/categories/},
  \texttt{GET /api/recipes/ingredients/},
  \texttt{GET /api/recipes/tags/};
  \item планы питания и списки покупок:
  \texttt{GET/POST /api/meal-plans/},
  \texttt{GET/PATCH/DELETE /api/meal-plans/\{id\}/},
  \texttt{POST /api/meal-plans/\{id\}/submit\_for\_moderation/},
  \texttt{POST /api/meal-plans/\{id\}/moderate/},
  \texttt{GET/POST /api/meal-plans/shopping-lists/},
  \texttt{POST /api/meal-plans/shopping-lists/sync-target/};
  \item пользовательские взаимодействия:
  \texttt{GET/POST /api/interactions/reviews/},
  \texttt{POST /api/interactions/reviews/\{id\}/moderate/},
  \texttt{GET/POST /api/interactions/complaints/},
  \texttt{POST /api/interactions/complaints/\{id\}/resolve/},
  \texttt{GET/POST /api/interactions/favorites/};
  \item уведомления и ассистент:
  \texttt{GET /api/notifications/},
  \texttt{POST /api/notifications/\{id\}/mark\_read/},
  \texttt{POST /api/assistant/chat/};
  \item служебные endpoint'ы:
  \texttt{GET /api/health/},
  \texttt{GET /api/schema/},
  \texttt{GET /api/docs/}.
\end{compactlist}

\section{Техническая реализация клиентской части}

Клиентская часть FoodPlanner реализована как одностраничное веб-приложение (SPA) на React + TypeScript. Выбор данного подхода обусловлен необходимостью обеспечить быстрый отклик интерфейса, типобезопасность при работе с данными API и удобство развития проекта при добавлении новых пользовательских сценариев.

В качестве UI-библиотеки используется Mantine, которая предоставляет готовые адаптивные компоненты и упрощает поддержку единого визуального стиля. Для управления состоянием применен Effector, позволяющий явно описывать события, эффекты и зависимости между состояниями.

Структура проекта клиентской части. Код приложения организован по модульному принципу, где каждый каталог отвечает за ограниченную область ответственности:
\begin{compactlist}
  \item \texttt{/Frontend/src/pages} --- страницы верхнего уровня (маршруты пользовательского интерфейса);
  \item \texttt{/Frontend/src/components} --- переиспользуемые компоненты каркаса и представления;
  \item \texttt{/Frontend/src/features} --- feature-модули со сценарием и состоянием (например, аутентификация, профиль, конструктор плана);
  \item \texttt{/Frontend/src/shared/api} --- централизованный слой HTTP-взаимодействия с backend и обработка ошибок;
  \item \texttt{/Frontend/src/shared/ui} --- общие UI-блоки (уведомления, виджет ассистента, состояния загрузки/пустых данных);
  \item \texttt{/Frontend/src/shared/model} --- общие клиентские модели и события;
  \item \texttt{/Frontend/src/types} --- доменные типы и DTO-представления данных.
\end{compactlist}

Структура проекта клиентской части представлена на рисунке~\ref{fig:image_19}.

\begin{figure}[H]
  \centering
  \IfFileExists{images/image_19.png}{
    \includegraphics[width=0.95\textwidth]{image_19.png}
  }{
    \fbox{\parbox{0.9\textwidth}{\centering Место для изображения структуры клиентской части\\
    Файл: \texttt{Report/images/image_19.png}}}
  }
  \caption{Структура проекта клиентской части}
  \label{fig:image_19}
\end{figure}

Слойная архитектура клиентского приложения включает:
\begin{compactlist}
  \item presentation layer --- React-страницы и компоненты визуализации;
  \item state layer --- модели Effector и локальные состояния интерфейса;
  \item data access layer --- API-функции, инкапсулирующие HTTP-запросы, CSRF-защиту, преобразование DTO и обработку ошибок.
\end{compactlist}

Разработка экранов и навигации. Экранная структура клиентской части соответствует пользовательским сценариям, рассмотренным в разделе 2.1: каталог рецептов, карточка рецепта, планы питания, конструктор, отзывы, профиль и авторизация. Для маршрутизации используется \texttt{react-router-dom}; базовый набор маршрутов приложения приведен в листинге~\ref{listing:frontend-routing}.

\begin{listing}[H]
\caption{Маршрутизация экранов клиентской части}
\label{listing:frontend-routing}
\begin{mdframed}
\begin{minted}[breaklines=true]{tsx}
<Routes>
  <Route path="/" element={<AboutPage />} />
  <Route path="/auth" element={<AuthPage />} />
  <Route path="/recipes" element={<RecipesPage />} />
  <Route path="/recipes/:recipeId" element={<RecipeDetailPage />} />
  <Route path="/meal-plans" element={<MealPlansPage />} />
  <Route path="/meal-plans/:planId" element={<MealPlanDetailPage />} />
  <Route path="/reviews" element={<ReviewsPage />} />
  <Route path="/planner" element={<PlannerPage />} />
  <Route path="/profile" element={<ProfilePage />} />
  <Route path="*" element={<FallbackCard message="Страница не найдена." />} />
</Routes>
\end{minted}
\end{mdframed}
\end{listing}

Управление состоянием в клиентской части реализовано через Effector. Для аутентификации используются события пользовательских действий, асинхронные эффекты и производные состояния статуса сессии. Пример модели авторизации приведен в листинге~\ref{listing:frontend-auth-model}.

\begin{listing}[H]
\caption{Модель состояния аутентификации на Effector}
\label{listing:frontend-auth-model}
\begin{mdframed}
\begin{minted}[breaklines=true]{ts}
export const checkSessionRequested = createEvent()
export const loginSubmitted = createEvent<LoginPayload>()
export const logoutRequested = createEvent()

export const checkSessionFx = createEffect(async () => getMe())
export const loginFx = createEffect(async (payload: LoginPayload) => login(payload))
export const logoutFx = createEffect(async () => logout())

export const $authStatus = createStore<AuthStatus>('checking')
  .on(checkSessionFx.done, () => 'auth')
  .on(checkSessionFx.fail, () => 'guest')
  .on(loginFx.done, () => 'auth')
  .on(logoutFx.done, () => 'guest')

sample({ clock: checkSessionRequested, target: checkSessionFx })
sample({ clock: loginSubmitted, target: loginFx })
sample({ clock: logoutRequested, target: logoutFx })
\end{minted}
\end{mdframed}
\end{listing}

Интеграция с серверной частью построена через централизованный API-слой с единым форматом вызовов \texttt{GET/POST/PATCH/DELETE}. Для state-changing запросов используется CSRF-токен, а также включен режим \texttt{credentials: include} для корректной работы с сессионной аутентификацией. Фрагмент клиентского API-слоя приведен в листинге~\ref{listing:frontend-api-csrf}.

\begin{listing}[H]
\caption{Клиентский API-слой с CSRF-защитой}
\label{listing:frontend-api-csrf}
\begin{mdframed}
\begin{minted}[breaklines=true]{ts}
export async function ensureCsrfCookie() {
  if (getCookie('csrftoken')) return
  await fetch(`${API_BASE_URL}/health/`, { credentials: 'include' })
}

async function apiRequest<T>(path: string, method: 'GET' | 'POST' | 'PATCH', body?: unknown) {
  if (method !== 'GET') await ensureCsrfCookie()
  const csrfToken = getCookie('csrftoken')

  const response = await fetch(`${API_BASE_URL}${path}`, {
    method,
    credentials: 'include',
    headers: method === 'GET' ? undefined : {
      'Content-Type': 'application/json',
      'X-CSRFToken': csrfToken,
    },
    body: body ? JSON.stringify(body) : undefined,
  })

  if (!response.ok) throw await toApiError(response)
  return response.status === 204 ? (undefined as T) : ((await response.json()) as T)
}
\end{minted}
\end{mdframed}
\end{listing}

Дополнительным пользовательским сценарием является работа с ИИ-ассистентом: в интерфейсе реализован плавающий виджет с открытием боковой панели, сохранением истории диалога, отправкой сообщений через backend-proxy endpoint и обработкой ошибок с выводом уведомлений.

\section{Вывод}

Во второй главе рассмотрены практические аспекты технической реализации системы FoodPlanner. Описаны интерфейсные решения, серверная архитектура, организация клиентского и серверного кода, а также ключевые принципы интеграции между слоями приложения. Представленные решения обеспечивают функциональную полноту пользовательских сценариев и создают устойчивую основу для дальнейшего развития платформы.